\section{Completeness of the irreducible locus}
\label{sec:complete2}

A natural worry is that \eqref{eq:family} is only a checked corner of a larger variety. It is not: it is
the entire irreducible locus for the principal spectrum.

\begin{theorem}[\sExact\ decomposition $+$ \sProved\ $\Fp$ classification]\label{thm:complete}
For the boundary $A$-spectrum $\{1,1,\omega,\omega^2\}$, the defining ideal of the figure-eight bundle
representations \eqref{eq:star}, stratified by the rank of the coupling block $Q$, has its primary
decomposition over $\Q(\omega)$ supported on the strata $\mathrm{rank}\,Q\in\{0,1,2\}$. Exactly one
component carries absolutely irreducible representations --- the four-parameter family
\eqref{eq:family} on $\mathrm{rank}\,Q=2$ --- and on it $[A,B]=-\mu^4$. Every other component is reducible
or vacuous ($\det t=0$). Consequently $L=-M^4$ holds \emph{unconditionally on the irreducible locus}.
\end{theorem}

\paragraph{\textmd{\textbf{What is proved, and how.}}} Two independent computations, with their status
stated honestly (this is the point the proof-status discipline exists for):
\begin{enumerate}[leftmargin=1.6em]
\item \emph{The component structure is symbolic} (\sExact). The gauge-rank stratification and the primary
  decomposition of each stratum's ideal are computed over $\Q(\omega)$ by Gr\"obner methods
  (minimal associated primes); this is an exact characteristic-zero statement, not a sample. The strata
  and their component counts are listed in Appendix~\ref{app:strata}.
\item \emph{Irreducibility is classified exactly over a finite field} (\sProved\ at the prime; see the
  caveat). On each component we apply the \emph{Burnside} criterion --- the algebra generated by
  $\rho(A),\rho(B)$ is the full matrix algebra $M_4$ iff $\rho$ is absolutely irreducible --- computed
  exactly over $\Fp$ ($p=1{,}000{,}003$) by span-closure. The result: only the $\mathrm{rank}\,Q=2$
  family returns full algebra (and all its samples satisfy $M^4=L$); every other component returns a
  proper subalgebra (reducible) or is vacuous.
\end{enumerate}

\paragraph{\textmd{\textbf{The Burnside--Schur caveat, and the validation.}}} The naive alternative ---
the Schur commutant being scalar --- is \emph{invalid here}: these representations are frequently
non-semisimple, and the commutant of a non-semisimple reducible representation can be scalar, so it
\emph{falsely certifies irreducibility}. (Indeed the verdict on one stratum flipped twice under scrutiny
before the Burnside test settled it.) We therefore use Burnside, and we \emph{validated} it against a
stratum proved reducible by hand --- the $\mathrm{rank}\,Q=0$ ($\det t=0$) locus, where a common invariant
subspace is exhibited symbolically (Appendix~\ref{app:strata}).

\paragraph{\textmd{\textbf{Honest scope.}}} What the two computations together establish in characteristic
zero is: the \emph{components} (symbolic), the irreducibility of the $\mathrm{rank}\,Q=2$ family
(a single Burnside witness lifts by Zariski-openness and rank semicontinuity), and the reducibility of the
$\det t=0$ stratum (symbolic invariant subspace). For the remaining non-principal components the
reducibility is established at exact $\Fp$ points; upgrading each to a characteristic-zero symbolic
invariant-subspace certificate is mechanical and is the one refinement we flag as in progress rather than
claim complete. We therefore badge Theorem~\ref{thm:complete} as \sExact\ in its component structure and
$\Fp$-exact in its irreducibility classification, \emph{not} as a single blanket \sProved; the
irreducible-locus conclusion is robust under this honest reading. Figure~\ref{fig:strat} summarises the
classification.

\begin{figure}[t]
  \centering
  \includegraphics[width=0.96\linewidth]{fig_strat}
  \caption{Gauge-rank strata of the bundle ideal, classified by the exact $\Fp$ Burnside test. Only the
  $\mathrm{rank}\,Q{=}2$ stratum (the family \eqref{eq:family}) carries irreducible representations, all
  satisfying $M^4=L$; the others are reducible or vacuous ($\det t=0$). The irreducible locus is the
  family. \sProved\ (at $p$)}
  \label{fig:strat}
\end{figure}
